% =====================================================================
%  Thème 2 — Angles associés — NIVEAU DIFFICILE
%  Corrigé
% =====================================================================
\documentclass[11pt,a4paper]{article}
\usepackage[utf8]{inputenc}
\usepackage[T1]{fontenc}
\usepackage{lmodern}
\usepackage[french]{babel}
\usepackage{amsmath,amssymb,mathtools}
\usepackage{icomma}
\usepackage{xcolor}
\usepackage{colortbl}
\usepackage{tikz}
\usepackage[most]{tcolorbox}
\usepackage{enumitem}
\usepackage{array}
\usepackage[a4paper,margin=2.2cm]{geometry}
\usetikzlibrary{arrows.meta,calc}

\definecolor{primary}{RGB}{214,51,108}
\definecolor{primarydark}{RGB}{140,20,64}
\definecolor{excol}{RGB}{0,135,90}

\DeclareMathOperator{\tg}{tg}
\DeclareMathOperator{\cotg}{cotg}
\newcommand{\degr}{^{\circ}}
\newcommand{\Z}{\mathbb{Z}}

\newcounter{exo}
\newcommand{\exo}[1]{\medskip\refstepcounter{exo}%
  \noindent{\large\bfseries\color{primarydark}Exercice \theexo}%
  \ \textcolor{primary}{\textbullet}\ \textbf{#1}\par\nopagebreak\smallskip}
\newcommand{\rep}{\textcolor{excol}{$\blacktriangleright$}\ }

\setlength{\parindent}{0pt}
\pagestyle{empty}

\begin{document}

\begin{center}
{\LARGE\bfseries\color{primary}Thème 2 — Angles associés}\\[2pt]
{\color{primarydark}\textsc{Niveau Difficile — Corrigé détaillé}}
\end{center}
\hrule height 1pt
\medskip

\exo{Un angle, deux chemins}
\begin{itemize}[leftmargin=1.4em,itemsep=2pt]
\item[\textbf{(a)}] \rep $\dfrac{5\pi}{6}=\pi-\dfrac{\pi}{6}$ \emph{(supplémentaires)}, donc
$\cos\dfrac{5\pi}{6}=-\cos\dfrac{\pi}{6}=-\dfrac{\sqrt3}{2}$.
\item[\textbf{(b)}] \rep $\dfrac{5\pi}{6}=\dfrac{\pi}{2}+\dfrac{\pi}{3}$ \emph{(anti-complémentaires)},
donc $\cos\dfrac{5\pi}{6}=-\sin\dfrac{\pi}{3}=-\dfrac{\sqrt3}{2}$.
\item[\textbf{(c)}] \rep Les deux méthodes donnent $-\dfrac{\sqrt3}{2}$ : elles \textbf{coïncident}.
C'est normal, $\pi-\dfrac{\pi}{6}$ et $\dfrac{\pi}{2}+\dfrac{\pi}{3}$ sont le même angle.
\item[\textbf{(d)}] \rep Supplémentaires : $\sin\dfrac{5\pi}{6}=\sin\dfrac{\pi}{6}=\dfrac12$.
Anti-complémentaires : $\sin\!\left(\dfrac{\pi}{2}+\dfrac{\pi}{3}\right)=\cos\dfrac{\pi}{3}=\dfrac12$.
Même résultat : $\sin\dfrac{5\pi}{6}=\dfrac12$.
\end{itemize}

\exo{Prouver une simplification}
\begin{itemize}[leftmargin=1.4em,itemsep=2pt]
\item[\textbf{(a)}] \rep $\sin(\pi-\alpha)=\sin\alpha$ \emph{(supplém.)} et $\sin(\pi+\alpha)=-\sin\alpha$
\emph{(anti-supplém.)}, donc $A=\sin\alpha-\sin\alpha=\mathbf{0}$.
\item[\textbf{(b)}] \rep $\cos(\pi+\alpha)=-\cos\alpha$ \emph{(anti-supplém.)} et $\cos(-\alpha)=\cos\alpha$
\emph{(opposés)}, donc $B=-\cos\alpha-\cos\alpha=\mathbf{-2\cos\alpha}$.
\item[\textbf{(c)}] \rep $\tg(\pi-\alpha)=-\tg\alpha$ \emph{(supplém.)} et $\tg(\pi+\alpha)=\tg\alpha$
\emph{(anti-supplém.)}, donc $C=-\tg\alpha+\tg\alpha=\mathbf{0}$.
\item[\textbf{(d)}] \rep $\sin\!\left(\dfrac{\pi}{2}+\alpha\right)=\cos\alpha$ \emph{(anti-complém.)} et
$\sin(\pi-\alpha)=\sin\alpha$ \emph{(supplém.)}, donc
$D=\cos^2\alpha+\sin^2\alpha=\mathbf{1}$ (identité fondamentale).
\end{itemize}

\exo{Quatre points associés}
\begin{itemize}[leftmargin=1.4em,itemsep=2pt]
\item[\textbf{(a)}] \rep En appliquant les familles à $M=(a\,;\,b)=(\cos\alpha\,;\,\sin\alpha)$ :
\[
N\ (\pi-\alpha)=(-a\,;\,b),\qquad
P\ (-\alpha)=(a\,;\,-b),\qquad
Q\ (\pi+\alpha)=(-a\,;\,-b).
\]
\item[\textbf{(b)}] \rep Même \textbf{ordonnée} (même $\sin$) : $M$ et $N$ (ordonnée $b$) ; $P$ et $Q$
(ordonnée $-b$). \ Même \textbf{abscisse} (même $\cos$) : $M$ et $P$ (abscisse $a$) ; $N$ et $Q$
(abscisse $-a$).
\item[\textbf{(c)}] \rep $M=(a\,;\,b)$ et $Q=(-a\,;\,-b)$ sont opposés coordonnée par coordonnée :
ils sont \textbf{symétriques par rapport au centre $O$} (anti-supplémentaires).
\item[\textbf{(d)}] \rep Pour $\alpha=\dfrac{\pi}{6}$ : $a=\cos\dfrac{\pi}{6}=\dfrac{\sqrt3}{2}$,
$b=\sin\dfrac{\pi}{6}=\dfrac12$. D'où
\[
M=\left(\tfrac{\sqrt3}{2}\,;\,\tfrac12\right),\
N=\left(-\tfrac{\sqrt3}{2}\,;\,\tfrac12\right),\
P=\left(\tfrac{\sqrt3}{2}\,;\,-\tfrac12\right),\
Q=\left(-\tfrac{\sqrt3}{2}\,;\,-\tfrac12\right).
\]
\end{itemize}

\exo{Toutes les écritures d'un même nombre}
On ramène chaque expression à une valeur remarquable.
\begin{itemize}[leftmargin=1.4em,itemsep=2pt]
\item[\textbf{(a)}] \rep $\sin\dfrac{5\pi}{6}=\sin\!\left(\pi-\dfrac{\pi}{6}\right)=\sin\dfrac{\pi}{6}=\dfrac12$
\emph{(supplém.)}. \textbf{Oui.}
\item[\textbf{(b)}] \rep $\cos\dfrac{2\pi}{3}=\cos\!\left(\pi-\dfrac{\pi}{3}\right)=-\cos\dfrac{\pi}{3}=-\dfrac12$
\emph{(supplém.)}. \textbf{Non} (vaut $-\tfrac12$).
\item[\textbf{(c)}] \rep $-\sin\!\left(-\dfrac{\pi}{6}\right)=-\left(-\sin\dfrac{\pi}{6}\right)=\sin\dfrac{\pi}{6}=\dfrac12$
\emph{(opposés)}. \textbf{Oui.}
\item[\textbf{(d)}] \rep $\sin\!\left(\dfrac{\pi}{2}+\dfrac{\pi}{3}\right)=\cos\dfrac{\pi}{3}=\dfrac12$
\emph{(anti-complém.)}. \textbf{Oui.}
\item[\textbf{(e)}] \rep $-\cos\dfrac{2\pi}{3}=-\left(-\dfrac12\right)=\dfrac12$ \emph{(supplém.)}. \textbf{Oui.}
\end{itemize}
Bilan : (a), (c), (d) et (e) valent $\dfrac12$ ; seule (b) vaut $-\dfrac12$.

\end{document}
